\documentclass[a4paper,11pt,openany]{ltjsreport}
\usepackage{lmodern}

% ============================================================
% フォント設定
% ============================================================
\usepackage{luatexja-fontspec}

\setmainfont{TeX Gyre Termes}
\setsansfont{TeX Gyre Heros}
\setmainjfont[
  BoldFont=Harano Aji Mincho Bold,
  ItalicFont=Harano Aji Mincho,
  BoldItalicFont=Harano Aji Mincho Bold
]{Harano Aji Mincho}

% ============================================================
% レイアウト
% ============================================================
\usepackage[top=25mm,bottom=25mm,left=25mm,right=25mm,bindingoffset=0mm]{geometry}
\usepackage{indentfirst}

% ============================================================
% 汎用パッケージ
% ============================================================
\usepackage{pdfpages}
\usepackage{siunitx}
\usepackage{enumitem}
\usepackage{graphicx}
\usepackage{amsmath}
\usepackage[normalem]{ulem}
\usepackage{url}
\usepackage{float}
\usepackage{amssymb}

% ============================================================
% 句読点変換設定 (Lua)
% ============================================================
\usepackage{luatexja}
\directlua{
  local function replace_punctuation(s)
    s = s:gsub("、", "，")
    s = s:gsub("。", "．")
    return s
  end
  luatexbase.add_to_callback("process_input_buffer", replace_punctuation, "replace_punctuation")
}

% ============================================================
% 表組み・図
% ============================================================
\usepackage{makecell}
\usepackage{hhline}
\usepackage{tabularray}
\usepackage{diagbox}
\usepackage{tikz}
\usepackage[framemethod=TikZ]{mdframed}
\makeatletter
\def\caption@documentclass{standard}
\makeatother
\usepackage{caption}
\usepackage{pifont}
\usepackage{varwidth}

% ============================================================
% 見出しスタイル・カラー・コードリスト
% ============================================================
\usepackage{titlesec}
\usepackage[dvipsnames]{xcolor}
\usepackage{listings}

% ============================================================
% 参考文献
% ============================================================
\usepackage[
  backend=biber,
  style=numeric,
  sorting=none,
  date=iso,
  urldate=iso,
  seconds=true,
  bibencoding=utf8
]{biblatex}
\addbibresource{ref/references.bib}

\DeclareFieldFormat*{title}{#1}
\DeclareFieldFormat*{journaltitle}{#1}
\DeclareFieldFormat*{booktitle}{#1}
\DeclareFieldFormat{url}{\url{#1}}
\DeclareDelimFormat{multinamedelim}{，}
\DeclareDelimFormat{finalnamedelim}{，}
\DefineBibliographyStrings{english}{
  urlseen = {参照},
}
\DeclareFieldFormat[article,inproceedings,misc]{title}{"#1"}
\DeclareFieldFormat[book,inbook]{title}{"#1"}

\DeclareFieldFormat{edition}{#1\addthinspace ed.}

% ============================================================
% キャプション日本語化
% ============================================================
\captionsetup[figure]{name=図}
\captionsetup[table]{name=表}
\captionsetup{labelsep=colon}

% ============================================================
% 図表・数式番号形式
% ----
% \chapter なし（図1 形式）→ 下記ブロックをそのまま使用
% \chapter あり（図1.1 形式）→ 下記ブロックを丸ごとコメントアウト
% ============================================================
\makeatletter
\@removefromreset{section}{chapter}
\@removefromreset{figure}{chapter}
\@removefromreset{table}{chapter}
\@removefromreset{equation}{chapter}
\makeatother
\renewcommand{\thesection}{\arabic{section}}
\renewcommand{\thefigure}{\arabic{figure}}
\renewcommand{\thetable}{\arabic{table}}
\renewcommand{\theequation}{\arabic{equation}}
% ============================================================

% ============================================================
% 見出しフォントサイズ・スペーシング
% ============================================================
\titleformat{\chapter}[block]
  {\fontsize{16pt}{18pt}\selectfont\bfseries}
  {第\thechapter 章\quad}
  {0pt}{}
\titlespacing*{\chapter}{0pt}{20pt}{15pt}

\titleformat{\section}
  {\normalsize\selectfont\bfseries}
  {(2-\thesection)\quad}
  {0pt}{}
\titlespacing*{\section}{0pt}{15pt}{3pt}

% subsection を、()付き小文字ローマン数字に変更。
\renewcommand{\thesubsection}{(\roman{subsection})}

\titleformat{\subsection}
  {\normalsize\selectfont\bfseries}
  {\thesubsection\quad}
  {0pt}{}
\titlespacing*{\subsection}{0pt}{10pt}{3pt}

\titleformat{\subsubsection}
  {\normalsize\bfseries}
  {\thesubsubsection\quad}
  {0pt}{}

\setcounter{secnumdepth}{3}

% ============================================================
% listings カラー定義
% ============================================================
\definecolor{codebg}{RGB}{245,245,245}
\definecolor{identifiercolor}{RGB}{7,54,66}
\definecolor{commentcolor}{RGB}{88,110,117}
\definecolor{framecolor}{RGB}{200,200,200}
\definecolor{keywordcolor}{RGB}{38,139,210}
\definecolor{builtincolor}{RGB}{133,153,0}
\definecolor{stringcolor}{RGB}{42,161,152}

% ============================================================
% listings スタイル定義
% ============================================================
\lstdefinestyle{codestyle}{
  backgroundcolor=\color{codebg},
  basicstyle=\linespread{1.1}\small\ttfamily\color{identifiercolor},
  frame=single,
  framerule=0.5pt,
  rulecolor=\color{framecolor},
  breaklines=true,
  showstringspaces=false,
  numbers=left,
  xleftmargin=2em,
  keywordstyle=\color{keywordcolor}\bfseries,
  stringstyle=\color{stringcolor},
  commentstyle=\color{commentcolor}\itshape,
  numberstyle=\normalsize\color{black},
}

\lstdefinestyle{python}{
  style=codestyle,
  language=Python,
  morekeywords=[2]{self,True,False,None,print,range,len,dict,list,set,zip,
    int,str,float,max,min,open,round,tuple,Exception,ValueError,
    FileNotFoundError,__name__},
  keywordstyle=[2]\color{builtincolor},
}

\lstdefinestyle{cpp}{
  style=codestyle,
  language=C++,
  morekeywords=[2]{
    nullptr,bool,true,false,
    uint8_t,uint16_t,uint32_t,
    int8_t,int16_t,int32_t,
    size_t,String,
    setup,loop,
    pinMode,digitalWrite,digitalRead,
    analogRead,analogWrite,
    delay,delayMicroseconds,
    millis,micros,
    HIGH,LOW,INPUT,OUTPUT,INPUT_PULLUP,
    Serial,begin,print,println,write,available,read,
    SoftwareSerial
  },
  keywordstyle=[2]\color{builtincolor},
}

\lstset{style=python}

% コードブロックの日本語化
\renewcommand{\lstlistingname}{コード}
\renewcommand{\lstlistlistingname}{コード一覧}

% ============================================================
% siunitx
% ============================================================
\sisetup{
  detect-all,
  mode=math
}

% ============================================================
% hyperref
% ============================================================
\usepackage{hyperref}
\hypersetup{
  colorlinks=true,
  linkcolor=black,
  urlcolor=black,
  citecolor=black,
  bookmarksnumbered=true,
  pdfborder={0 0 0},
  pdfencoding=auto,
  pdftitle={OS 前期末レポート1},
  pdfauthor={中山 将吾}
}

\begin{document}

\setcounter{page}{1}

% ============================================================
% 表紙（どちらか一方を有効化）
% ============================================================

% -- 指定表紙PDFを使う場合（通常はこちら） --
% \includepdf{./cover.pdf}

% -- 自作表紙を使う場合はコメントを外す --
\begin{titlepage}
	\vspace*{\stretch{1.0}}
	\begin{flushright}
		\large{提出日: 2026年7月6日}
	\end{flushright}
	\vspace*{\stretch{2.3}}
	\begin{center}
		{\huge オペレーティングシステム} \\
		{\vspace{1cm}}
		{\huge レポート課題3}
	\end{center}
	\vspace*{\stretch{3.0}}
	\begin{center}
		{\large 所属: 電子情報工学科・4年} \\
		\vspace{0.5cm}
		{\large 学籍番号: i231331} \\
		\vspace{0.5cm}
		{\large 出席番号: 33番} \\
		\vspace{0.5cm}
		{\large 氏名: 中山 将吾}
	\end{center}
	\vspace*{\stretch{1.0}}
\end{titlepage}

\section*{問題}
以下の仕様のディスクから$10\mathrm{MB}$のデータをランダムアクセスで読み出すのにかかる最悪時間と平均時間の両方を求めよ。
加えて、シーケンシャルアクセスした場合と比較し、ランダムアクセスがどれほど非効率かを倍率で答えよ。
ブロックサイズを$4\mathrm{kB}$とし、1ブロックずつ個別に読み出すものとする。
\begin{itemize}
	\item ディスク回転数: $6000\mathrm{rpm}$
	\item 平均シーク時間: $\SI{20}{\ms}$
	\item ディスクのデータ転送速度(シーケンシャルアクセス時): $40\mathrm{MB/s}$
\end{itemize}

\subsection*{解答}
以下の手順で求める。
\begin{enumerate}[label=(\arabic*)]
	\item ランダムアクセスの回数を求める。
	\item 1回のランダムアクセスに要するアクセス時間を求める。
	\item これらを踏まえて解を求めよ。
\end{enumerate}

\begin{enumerate}[label=(\arabic*)]
	\item \textbf{ランダムアクセスの回数}\\
	      データサイズが$10\mathrm{MB}$であり、$1\mathrm{MB} = 1000\mathrm{kB}$より、
	      \begin{align*}
		      10\mathrm{MB} = 10000\mathrm{kB}
	      \end{align*}
	      である。$1$ブロックは$4\mathrm{kB}$であるため、アクセス回数は
	      \begin{align*}
		      \frac{10000\mathrm{kB}}{4\mathrm{kB}} = 2500\mathrm{回}
	      \end{align*}
	      となる。

	\item \textbf{1回のランダムアクセスに要するアクセス時間}\\
	      1回のアクセス時間は、シーク時間、回転待ち時間、データ転送時間の合計である。
	      \begin{enumerate}[label=(\roman*)]
		      \item \textbf{シーク時間}\\
		            問題文より、シーク時間は$\SI{20}{\ms}$である。

		      \item \textbf{回転待ち時間}\\
		            ディスク回転数は$6000\mathrm{rpm}$であるため、1回転にかかる時間は、
		            \begin{align*}
			            \frac{\SI{60}{\s}}{6000\mathrm{rpm}} & = \SI{0.01}{\s} \\
			                                                 & = \SI{10}{\ms}
		            \end{align*}
		            となる。
		            平均回転待ち時間は1回転の半分であるため$\SI{5}{\ms}$であり、最悪回転待ち時間は1回転分の$\SI{10}{\ms}$となる。

		      \item \textbf{データ転送時間}\\
		            1トラックあたりの記憶容量が$20\mathrm{kB}$であり、この読み出しに1回転($\SI{10}{\ms}$)かかるとみなすと、$4\mathrm{kB}$のブロックを読み出す時間は、
		            \begin{align*}
			            \SI{10}{\ms} \times \left(\frac{4\mathrm{kB}}{20\mathrm{kB}}\right) = \SI{2}{\ms}
		            \end{align*}
	      \end{enumerate}

	      これらより、1回あたりのアクセス時間は
	      \begin{itemize}
		      \item \textbf{平均アクセス時間}
		            \begin{align*}
			            \SI{20}{\ms} + \SI{5}{\ms} + \SI{2}{\ms} = \SI{27}{\ms}
		            \end{align*}
		      \item \textbf{最悪アクセス時間}
		            \begin{align*}
			            \SI{20}{\ms} + \SI{10}{\ms} + \SI{2}{\ms} = \SI{32}{\ms}
		            \end{align*}
	      \end{itemize}

	\item \textbf{これらを踏まえた解}\\
	      (1)で求めたアクセス回数と(2)で求めた1回あたりのアクセス時間を乗算する。
	      \begin{itemize}
		      \item \textbf{ランダムアクセス平均時間}
		            \begin{align*}
			            \SI{27}{\ms} \times 2500\mathrm{回} & = \SI{67500}{\ms} \\
			                                               & = \SI{67.5}{\s}
		            \end{align*}

		      \item \textbf{ランダムアクセス最悪時間}
		            \begin{align*}
			            \SI{32}{\ms} \times 2500\mathrm{回} & = \SI{80000}{\ms} \\
			                                               & = \SI{80.0}{\s}
		            \end{align*}
	      \end{itemize}
	      シーケンシャルアクセスした場合の時間は、データ量とデータ転送速度から
	      \begin{align*}
		      \frac{10\mathrm{MB}}{40\mathrm{MB/s}} = \SI{0.25}{\s}
	      \end{align*}
	      ランダムアクセスがどれほど非効率かを倍率で計算すると、
	      \begin{itemize}
		      \item \textbf{平均時間の倍率}
		            \begin{align*}
			            \frac{\SI{67.5}{\s}}{\SI{0.25}{\s}} = 270\mathrm{倍}
		            \end{align*}

		      \item \textbf{最悪時間の倍率}
		            \begin{align*}
			            \frac{\SI{80.0}{\s}}{\SI{0.25}{\s}} = 320\mathrm{倍}
		            \end{align*}
	      \end{itemize}

	      よって、ランダムアクセスにかかる時間と非効率の倍率は以下の通りである。
	      \begin{itemize}
		      \item 平均時間: \SI{67.5}{\s} ($270$倍)
		      \item 最悪時間: \SI{80.0}{\s} ($320$倍)
	      \end{itemize}
\end{enumerate}

% \printbibliography[title=参考文献]

\end{document}
